% Trellis-Coded Quantization for KV Cache Compression

\documentclass{article}
\usepackage[utf8]{inputenc}
\usepackage[T1]{fontenc}
\usepackage{libertine}
\usepackage{amsmath,amsthm}
\usepackage[libertine]{newtxmath}
\usepackage[margin=1.25in]{geometry}
\usepackage{graphicx}
\usepackage{booktabs}
\usepackage[hidelinks]{hyperref}
\usepackage{algorithm}
\usepackage{algorithmic}
\usepackage{xcolor}
\usepackage{titlesec}
\usepackage{tcolorbox}

% --- Theme colors ---
\definecolor{themecolor}{RGB}{0, 0, 0}
\definecolor{themebg}{RGB}{245, 245, 248}
\definecolor{themeborder}{RGB}{190, 190, 200}

% Colored section headers
\titleformat{\section}
  {\Large\bfseries\color{themecolor}}{\thesection}{1em}{}
\titleformat{\subsection}
  {\large\bfseries\color{themecolor}}{\thesubsection}{1em}{}
\titleformat{\subsubsection}
  {\normalsize\bfseries\color{themecolor}}{\thesubsubsection}{1em}{}

\newtheorem{theorem}{Theorem}
\newtheorem{proposition}{Proposition}

% Add breathing room around floats
\setlength{\textfloatsep}{18pt plus 4pt minus 4pt}
\setlength{\intextsep}{14pt plus 4pt minus 4pt}
\setlength{\floatsep}{14pt plus 4pt minus 4pt}
\setlength{\abovecaptionskip}{10pt}
\setlength{\belowcaptionskip}{4pt}
\newcommand{\R}{\mathbb{R}}
\newcommand{\E}{\mathbb{E}}
\newcommand{\Var}{\mathrm{Var}}

\begin{document}

% --- Custom left-aligned title ---
\noindent
{\LARGE\bfseries Closing the Gap:}\\[4pt]
{\LARGE\bfseries Trellis-Coded Quantization for KV Cache at 2--3 Bits}\\[12pt]
{\large buun$^{\dagger}$, Claude (Anthropic)}\\[4pt]
\rule{\textwidth}{0.4pt}
\vspace{1.5em}

% ============================================================
\begin{tcolorbox}[
  colback=themebg,
  colframe=themeborder,
  arc=4pt,
  boxrule=0.6pt,
  left=12pt, right=12pt, top=10pt, bottom=8pt,
  fontupper=\small
]
\centerline{\normalsize\bfseries\color{themecolor} Abstract}
\vspace{0.4em}

We present the first application of trellis-coded quantization (TCQ) to
key-value cache compression in large language model inference. By combining
Fast Walsh-Hadamard rotation with Viterbi-optimal trellis encoding, our method
achieves 10--44\% KL-divergence reduction over scalar quantization at 2--3 bits
per value, with $O(1)$ parallel decode. At 3.25 bits per value,
our method produces \emph{lower} perplexity than uncompressed FP16 KV cache
due to a beneficial norm-scaling effect. At 2.25 bits per value, TCQ
significantly improves 2-bit quantization, closing the gap with 3-bit scalar methods.
We further identify a context-dependent codebook quality phenomenon: the
relationship between codebook MSE and downstream model quality inverts between
short and long contexts, explained by CLT averaging of quantization error in the
attention mechanism.

\vspace{0.5em}
\footnotesize
$^{\dagger}$\textbf{Correspondence:} \url{https://x.com/spiritbuun}\\
\textbf{Code:} {\small\url{https://github.com/spiritbuun/llama-cpp-turboquant-cuda}}
\end{tcolorbox}

% ============================================================
\section{Introduction}
\label{sec:intro}

Long-context inference is central to modern large language model (LLM)
applications — from document analysis and code generation to multi-turn
dialogue. The key-value (KV) cache, which stores intermediate attention states
to avoid redundant computation, grows linearly with sequence length and
dominates GPU memory consumption during inference. For a 27-billion-parameter
model at 128K context, the KV cache alone exceeds 40\,GB in FP16, often
surpassing the model weights themselves.

Quantizing the KV cache to low bit widths (2--4 bits per value) is the most
direct approach to this memory bottleneck. Recent work has made significant
progress using scalar quantization~\cite{kivi2024,kvquant2024} and online
vector quantization with random rotation~\cite{turboquant2026}. However, all
existing methods quantize each element or vector independently: they do not
exploit sequential structure within the KV vectors to achieve compression
beyond what independent coding permits.

In the signal processing literature, \emph{trellis-coded quantization}
(TCQ)~\cite{marcellin1990tcq} has long been known to close most of the 1.53\,dB
gap between scalar and optimal vector quantization for Gaussian sources. TCQ
constrains the sequence of quantization indices to follow a finite-state
trellis, enabling a much larger effective codebook at the same bit rate. The
Viterbi algorithm finds the globally optimal encoding, and — crucially for GPU
inference — the bitshift trellis structure allows each element to be decoded
independently in $O(1)$ via a sliding bit window.

Despite TCQ's theoretical appeal, it has not been applied to KV cache
compression. The only prior application of TCQ to neural networks is QTIP's
use for weight quantization~\cite{qtip2024}, where the encoding cost is
amortized at model load time. For KV cache, encoding must happen online
during inference, demanding a lean trellis and efficient CUDA implementation.

In this work, we combine Fast Walsh-Hadamard rotation — which transforms KV
vectors into an i.i.d.\ Gaussian distribution amenable to analytical codebook
design — with Viterbi-optimal trellis encoding using a 512-state bitshift
trellis. Our method stores only 3.25 bits per value (3-bit) or 2.25 bits per
value (2-bit) while achieving substantial quality improvements over scalar
quantization at the same bit rate.

\textbf{Contributions.} We make four contributions:

\begin{enumerate}
\item \textbf{TCQ for KV cache.} We present the first application of
  trellis-coded quantization to key-value cache compression, achieving
  10--44\% KL-divergence reduction over scalar quantization at 2--3 bits per
  value (Section~\ref{sec:method}). At 3.25 bpv, our method produces lower
  perplexity than uncompressed FP16 due to a beneficial norm-scaling effect.
  At 2.25 bpv, TCQ substantially improves 2-bit quantization,
  reducing the PPL gap from $+$9\% to $+$1.2\% relative to FP16.

\item \textbf{Context-dependent codebook quality.} We identify and explain a
  phenomenon where the relationship between codebook MSE and model quality
  inverts between short and long contexts
  (Section~\ref{sec:context_codebook}). At short contexts, codebook
  \emph{regularity} determines quality; at long contexts, \emph{MSE}
  dominates via CLT averaging of quantization error in the attention
  mechanism. No single codebook is optimal across all context lengths.

\item \textbf{MSE--quality divergence.} We demonstrate that every aggregate
  error metric we tested (13 statistics including covariance structure,
  kurtosis, tail behavior, and simulated attention error) improves
  monotonically with codebook training iterations, yet model perplexity
  degrades beyond 10 iterations (Section~\ref{sec:mse_divergence}). This
  provides empirical evidence for the theoretical prediction that MSE is the
  wrong objective for quantization in attention~\cite{ordentlich2024matmul}.

\item \textbf{Query anisotropy.} We measure extreme anisotropy in the rotated
  query space (effective rank $\sim$5 out of 128 dimensions), implying that
  96\% of quantization error is invisible to the attention mechanism
  (Section~\ref{sec:anisotropy}). Current TCQ errors are isotropic in the
  query eigenbasis, suggesting a 25$\times$ theoretical ceiling for
  query-aware codebook design.
\end{enumerate}

\noindent Our implementation is open-source, integrated into llama.cpp with full CUDA
support for both encoding and flash-attention decoding.

% ============================================================
\section{Background and Related Work}
\label{sec:background}

We situate our work at the intersection of three lines of research: KV cache
compression for efficient LLM inference, trellis-coded quantization from the
signal processing literature, and random rotation techniques for quantization.

\subsection{KV Cache Compression}
\label{sec:bg_kvcache}

The key-value cache stores intermediate attention states across all layers and
grows linearly with sequence length, consuming the majority of GPU memory
during long-context inference. For a model with $L$ layers, $h$ attention
heads, head dimension $d$, and sequence length $n$, the KV cache requires
$2Lhdn$ elements in FP16 — over 40\,GB for a 27B-parameter model at 128K
context.

\textbf{Scalar quantization.} The simplest approach quantizes each KV element
independently. llama.cpp's built-in types (q4\_0, q8\_0) use per-block
quantization with block sizes of 32, achieving 4--8 bits per value with
modest quality loss~\cite{llama_cpp}. KIVI~\cite{kivi2024} applies per-channel
quantization with 2-bit keys and 2-bit values, using residual compensation
to maintain quality. KVQuant~\cite{kvquant2024} adds per-channel sensitivity
weighting and non-uniform quantization for sub-4-bit operation.

\textbf{Vector quantization.} TurboQuant~\cite{turboquant2026} introduced
online vector quantization for KV cache, using random rotation to
standardize the distribution before applying analytically-computed Lloyd-Max
codebooks. This eliminates the need for calibration data and achieves
near-optimal distortion at 3--4 bits per value. VecInfer~\cite{vecinfer2025}
combines smooth transforms with Hadamard rotation and VQ, with fused CUDA
kernels achieving 2.7$\times$ speedup. NSNQuant~\cite{nsnquant2025} uses triple
normalization to align distributions to standard normal, enabling calibration-free
codebooks — an observation that aligns with our post-FWHT Gaussian finding.

\textbf{Transform coding.} KVTC~\cite{kvtc2025} applies classical transform
coding (PCA decorrelation + adaptive quantization + entropy coding) to achieve
up to 20$\times$ compression. Their PCA-based decorrelation is data-dependent
and requires calibration, in contrast to the data-independent FWHT rotation
used in our work and TurboQuant.

\textbf{Bit allocation.}
``More Keys, Less Values''~\cite{morekeyslessvalues2025} proves theoretically
that keys have higher spectral norms than values, motivating asymmetric
bit budgets (e.g., 4-bit keys with 2-bit values) — a finding we validate
empirically in Section~\ref{sec:results}.

\textbf{Our position.} All of the above use either scalar quantization or
independent vector quantization — each element or vector is quantized without
considering its neighbors in the sequence. We introduce \emph{trellis-coded}
quantization, which jointly optimizes the quantization of 128 consecutive
elements via the Viterbi algorithm, achieving a larger effective codebook at
the same bit rate. To our knowledge, no prior work has applied TCQ to KV cache
compression.

\subsection{Trellis-Coded Quantization}
\label{sec:bg_tcq}

Trellis-coded quantization was introduced by Marcellin and
Fischer~\cite{marcellin1990tcq} as the source-coding dual of Ungerboeck's
trellis-coded modulation (TCM). The key idea: by constraining the sequence of
quantization indices to follow a trellis (finite-state machine), one can use a
larger codebook than the bit rate would otherwise permit, with the trellis
structure ensuring decodability.

\textbf{Coding gain.} For an i.i.d.\ Gaussian source, a scalar quantizer at
rate $R$ bits/sample incurs a distortion penalty of 1.53\,dB relative to the
Shannon rate-distortion bound $D(R) = \sigma^2 \cdot 2^{-2R}$. This gap
is fundamental to one-dimensional quantization~\cite{gray1998quantization}.
TCQ recovers most of this gap: an 8-state TCQ achieves $\sim$0.5\,dB from
$R(D)$, and a 256-state TCQ achieves $\sim$0.2\,dB~\cite{fischer1992ectcq}.
The gain comes from the trellis effectively creating a high-dimensional
quantizer from scalar operations — the ``effective dimension'' grows with
trellis memory~\cite{gersho1979,lookabaugh1989}.

\textbf{Structure.} A TCQ with $2^L$ states and rate $k$ bits/element has a
codebook of $2^L$ reconstruction values, organized into $2^k$ cosets via set
partitioning. At each time step, the trellis state determines which cosets are
reachable; within each reachable coset, the nearest codeword is selected. The
Viterbi algorithm finds the globally optimal path in $O(2^L \cdot n)$ time.

\textbf{Universal TCQ.} Kasner and Marcellin~\cite{kasner1999utcq} showed
that TCQ with \emph{uniform thresholds} (no training) achieves performance
``comparable with fully optimized entropy-constrained TCQ'' at most encoding
rates. This remarkable result implies that the coding gain comes primarily from
the trellis structure and set partitioning, not from codeword optimization.
However, as we show in Section~\ref{sec:context_codebook}, this universality
breaks down at very low rates (2-bit) where codebook training is essential.

\textbf{TCQ for neural network quantization.} QTIP~\cite{qtip2024} recently
applied TCQ to LLM weight quantization with impressive results at 2--4 bits
per weight. QTIP uses a bitshift trellis with $2^{16}$ states, hash-based
codebook mapping, and Hessian-weighted distortion (LDLQ). Our work differs
in three respects: (1) we quantize activations (KV cache) rather than weights,
requiring online encoding at inference time; (2) we use a much smaller trellis
($2^9 = 512$ states) to keep encode latency manageable; (3) we exploit the
bitshift trellis's sliding-window property for $O(1)$ parallel GPU decode,
which is not needed for weight dequantization (done once at load time).

\subsection{Random Rotation for Quantization}
\label{sec:bg_rotation}

Channel-wise outliers in LLM activations cause severe quantization
degradation: a few high-magnitude channels dominate the dynamic range,
starving other channels of precision. Random orthogonal rotation
redistributes energy uniformly across channels, eliminating outliers at the
cost of one matrix-vector multiply.

\textbf{Theory.} For a unit-norm vector $\mathbf{x} \in \R^d$, the
coordinates of $R\mathbf{x}$ (where $R$ is drawn uniformly from the
orthogonal group) are approximately i.i.d.\ $\mathcal{N}(0, 1/d)$ for large
$d$, by concentration of measure on the sphere. The Fast Walsh-Hadamard
Transform with random sign flips provides an efficient $O(d \log d)$
approximation to a random rotation~\cite{ailon2009fast}.

\textbf{Applications.} QuIP\#~\cite{quipsharp2024} combines random Hadamard
rotation with E8 lattice codebooks for weight quantization, demonstrating
that rotation + structured codebook outperforms trained codebooks despite the
latter achieving lower MSE on training data. This finding, grounded in the
Zamir-Feder theorem that lattice quantization noise is approximately
white~\cite{zamir1996lattice}, parallels our observation that codebook
regularity matters more than MSE at short contexts
(Section~\ref{sec:context_codebook}).

TurboQuant~\cite{turboquant2026} applies FWHT rotation to KV cache
quantization and proves that the resulting distribution admits analytically
optimal scalar quantizers within 2.72$\times$ of the Shannon bound — no
training data required. RotateKV~\cite{rotatekv2025} extends this with
channel reordering before FWHT to further reduce outlier impact.
MixQuant~\cite{mixquant2026} provides the first systematic analysis of
outlier suppression under block Hadamard rotations, showing that outliers are
minimized when pre-rotation $L_1$ mass is evenly distributed.

\textbf{Our contribution.} We combine FWHT rotation (which makes the data
amenable to analytical codebook design) with TCQ (which provides coding gain
beyond scalar quantization). The rotation converts the KV quantization problem
from an intractable high-dimensional distribution to a well-understood i.i.d.\
Gaussian, for which TCQ's theoretical properties are fully
characterized~\cite{marcellin1990tcq}. This combination — rotation for
distribution standardization, trellis coding for compression efficiency — has
not been previously explored for KV cache or, to our knowledge, for any
online activation quantization setting.

% ============================================================
\section{Method}
\label{sec:method}

We describe our approach to compressing the KV cache using trellis-coded
quantization. The pipeline processes each 128-element group of a key or value
vector through five stages: (1) Walsh-Hadamard rotation, (2) norm extraction,
(3) Viterbi trellis encoding, (4) bitpacked storage, and (5) sliding-window
decode during attention. Figure~\ref{fig:pipeline} illustrates the full
pipeline.

% pipeline figure would go here
% \begin{figure}[t]
% \centering
% \includegraphics[width=\columnwidth]{figures/pipeline.pdf}
% \caption{TCQ KV cache pipeline.}
% \label{fig:pipeline}
% \end{figure}

\subsection{Walsh-Hadamard Rotation}
\label{sec:rotation}

Raw KV vectors exhibit highly non-uniform channel distributions: a few channels
carry most of the energy while others are near-zero. Directly quantizing such
vectors wastes bits on low-energy channels and underserves high-energy ones.

Following TurboQuant~\cite{turboquant2026}, we apply the Fast Walsh-Hadamard
Transform (FWHT) with random sign flips to each 128-element group before
quantization. For a group $\mathbf{x} \in \R^{128}$, we compute:

\begin{equation}
\tilde{\mathbf{x}} = \frac{1}{\sqrt{128}} H_{128} \, S \, \mathbf{x}
\end{equation}

where $H_{128}$ is the $128 \times 128$ Hadamard matrix (applied in
$O(n \log n)$ via the butterfly algorithm) and $S = \mathrm{diag}(s_1, \ldots,
s_{128})$ with $s_i \in \{-1, +1\}$ drawn from a fixed pseudorandom sequence.

After rotation, each coordinate of $\tilde{\mathbf{x}}$ is approximately i.i.d.\
Gaussian with variance $\|\mathbf{x}\|^2 / 128$. This is a consequence of the
concentration of measure on the sphere: for high-dimensional vectors, random
orthogonal projections yield near-Gaussian marginals.

This transformation has three critical effects:
\begin{enumerate}
\item \textbf{Uniformization.} All channels have equal variance, eliminating the
  need for per-channel scaling or mixed-precision allocation.
\item \textbf{Known distribution.} The post-rotation distribution is analytically
  characterized (Gaussian), enabling optimal codebook design without calibration
  data.
\item \textbf{Decorrelation.} Inter-channel correlations are destroyed, making
  the i.i.d.\ assumption required by TCQ theory valid in practice.
\end{enumerate}

The rotation is applied once at KV cache write time and inverted once at
attention compute time, adding negligible overhead to both paths.

\subsection{Norm Extraction and Scaling}
\label{sec:norm}

Before trellis encoding, we extract the $L_2$ norm of each 128-element group:

\begin{equation}
n = \|\tilde{\mathbf{x}}\|, \qquad \hat{\mathbf{x}} = \tilde{\mathbf{x}} / n
\end{equation}

The norm $n$ is stored as a single FP16 value shared across 128 elements
(0.125 bits/element overhead). The normalized vector $\hat{\mathbf{x}}$ has
unit norm with each coordinate distributed as $\mathcal{N}(0, 1/128)$.

During reconstruction, the decoded vector is rescaled by $\alpha \cdot n$,
where $\alpha$ is a learned scaling factor. We find that the optimal $\alpha$
differs for keys and values:

\begin{equation}
\hat{\mathbf{x}}_{\text{recon}} =
\begin{cases}
\alpha_K \cdot n \cdot \mathrm{decode}(\text{bits}) & \text{for keys} \\
\alpha_V \cdot n \cdot \mathrm{decode}(\text{bits}) & \text{for values}
\end{cases}
\end{equation}

with $\alpha_K = 1.0$ and $\alpha_V$ set based on the target context depth.
The optimal $\alpha_V$ varies with context length
(Table~\ref{tab:alpha}): at 3-bit, it decreases from 1.04 at 2K to 1.00
at 32K; at 2-bit, it increases from 1.04 to $\sim$1.10.

Alpha can be applied at either \emph{encode time} (baked into the fp16 norm
before storage) or \emph{decode time} (multiplied during dequantization in
the attention kernel). Decode-time application is the preferred deployment
path: it enables \textbf{context-adaptive alpha} --- the optimal scaling can
be adjusted based on the current context depth without re-encoding the KV
cache. At contexts of 8K and above, decode-time alpha actually produces
lower KLD than encode-time at matched alpha values (Section~\ref{sec:results}).

The asymmetry arises because key and value errors propagate differently through
attention. Key errors enter inside the softmax exponential and are nonlinearly
amplified~\cite{asymkv2024}, while value errors propagate linearly through the
weighted average. The slight upward scaling of value norms provides a beneficial
regularization effect: it sharpens the attention distribution, reducing the
influence of noisy low-weight tokens. At 3.25 bits per value, this scaling is
sufficient for the quantized model to achieve \emph{lower perplexity than
uncompressed FP16} (Section~\ref{sec:results}).

\subsection{Trellis-Coded Quantization}
\label{sec:tcq}

\subsubsection{Trellis Structure}

We use a bitshift trellis with $2^L = 512$ states, rate $k = 3$ bits per
element, and scalar quantization ($V = 1$). The trellis is a directed graph
where state $i$ has an edge to state $j$ if and only if:

\begin{equation}
j = (i \gg k) \,|\, (c \ll (L - k))
\quad \text{for some } c \in \{0, 1, \ldots, 2^k - 1\}
\label{eq:transition}
\end{equation}

That is, the next state is formed by right-shifting the current state by $k$
bits and inserting $k$ new bits at the top. Each state has $2^k = 8$ outgoing
edges and $2^k = 8$ incoming edges.

Each state $s \in \{0, \ldots, 2^L - 1\}$ is associated with a reconstruction
value $C[s] \in \R$ from a trained codebook. The codebook is organized into
$2^k = 8$ \emph{cosets}, where coset $c$ contains all states with the same
upper $k$ bits. Set partitioning ensures that states within the same coset are
well-separated in reconstruction value, analogous to Ungerboeck's
trellis-coded modulation~\cite{marcellin1990tcq}.

\subsubsection{Viterbi Encoding}

Given a normalized vector $\hat{\mathbf{x}} = (x_1, \ldots, x_{128})$, we find
the trellis walk that minimizes total squared error:

\begin{equation}
\min_{s_1, \ldots, s_{128}} \sum_{t=1}^{128} (x_t - C[s_t])^2
\quad \text{s.t. } (s_t, s_{t+1}) \text{ is a valid edge } \forall t
\label{eq:viterbi_obj}
\end{equation}

This is solved exactly by the Viterbi algorithm. At each time step $t$, we
maintain a cost vector over all $2^L$ states and update via:

\begin{equation}
V_t(j) = \min_{i : (i,j) \in E} \left[ V_{t-1}(i) + (x_t - C[j])^2 \right]
\end{equation}

After processing all 128 elements, we backtrack from the minimum-cost final
state to recover the optimal state sequence.

\textbf{Free initialization.} Unlike classic TCQ which constrains the initial
state, we allow all $2^L$ states as valid starting points with equal cost. This
yields 37.6\% additional MSE reduction at 3-bit by removing an artificial
constraint on the first few elements.

\textbf{Complexity.} Encoding is $O(2^L \times n)$ per block, where $n = 128$.
With $L = 9$, this is $128 \times 512 = 65{,}536$ state evaluations per block.
The Viterbi iterations are \emph{inherently sequential} along the time axis —
this is the fundamental cost of trellis coding. In our CUDA implementation,
encode uses 512 threads with barrier-synchronized iterations and runs entirely
in shared memory.

\subsubsection{Bitpacked Storage}

The trellis walk is stored as a packed bitstream. The initial state requires $L$
bits; each subsequent step contributes only $k$ new bits (the edge label $c$
from Equation~\ref{eq:transition}). For 128 elements at $k = 3$ bits:

\begin{equation}
\text{Total} = L + 127 \times k = 9 + 381 = 390 \text{ bits}
\end{equation}

Combined with the FP16 norm (16 bits shared across 128 elements), this gives a
storage cost of $(390 + 16) / 128 = 3.17$ bits per value for 3-bit TCQ. In
practice, byte-alignment padding brings this to 3.25 bpv. At 2-bit ($k = 2$):
$(9 + 127 \times 2 + 16) / 128 = 2.23$ bpv, padded to 2.25 bpv.

\subsubsection{Sliding-Window Decode}
\label{sec:decode}

The critical property enabling GPU-efficient inference: \textbf{each element can
be decoded independently in $O(1)$}.

To reconstruct element $t$, we read $L$ consecutive bits starting at position
$t \times k$ in the bitstream. These $L$ bits \emph{are} the state index $s_t$,
because the bitshift transition rule ensures that consecutive states share
$L - k$ bits:

\begin{equation}
s_t = \mathrm{bits}[t \cdot k \,:\, t \cdot k + L]
\end{equation}

The reconstructed value is then $C[s_t] \times \alpha \times n$ — a single
codebook lookup and two multiplications. The 512-entry codebook (2\,KB) fits
comfortably in GPU shared memory, enabling coalesced access with no global
memory traffic during attention computation.

This $O(1)$ decode property is fundamental: while encoding requires sequential
Viterbi (done once at KV cache write time), \emph{decoding is fully parallel}
and sits on the latency-critical attention path. Each token's KV cache is
decoded independently across all heads and layers.

\subsection{Effective Codebook Size}
\label{sec:effective_codebook}

A key advantage of TCQ over scalar quantization is the \emph{effective codebook
expansion}. A scalar $k$-bit quantizer has $2^k$ reconstruction levels. TCQ
with $2^L$ states has $2^L$ reconstruction levels while storing only $k$ bits
per element.

At 3-bit with $L = 9$: the scalar quantizer chooses among 8 levels; TCQ chooses
among 512 levels, constrained by the trellis. This $64\times$ expansion of the
effective codebook is what provides TCQ's coding gain — the ability to represent
the source with finer granularity than the bit rate would otherwise permit.

The coding gain of an $N$-state TCQ for a Gaussian source is well-characterized
in the literature~\cite{marcellin1990tcq,fischer1992ectcq}: 8-state TCQ
achieves approximately 0.5~dB from the rate-distortion bound $R(D)$; 256-state
TCQ achieves approximately 0.2~dB. Our 512-state trellis operates in the regime
of diminishing returns, capturing the majority of available coding gain.


% ============================================================
\section{Codebook Design and Theory}
\label{sec:codebook}

The design of TCQ codebooks for KV cache quantization reveals several phenomena
not predicted by classical quantization theory. This section presents our
findings on codebook training, the divergence between MSE and model quality, and
the context-length dependence of optimal codebook selection.

\subsection{Codebook Training via Generalized Lloyd Algorithm}
\label{sec:gla}

Our codebook training follows the Generalized Lloyd Algorithm (GLA) applied to
the TCQ trellis~\cite{sabin1986gla}. The procedure alternates two steps:

\begin{enumerate}
\item \textbf{Encode step.} Given the current codebook $C$, run Viterbi
  encoding on training data to find the optimal state sequence for each training
  vector.
\item \textbf{Centroid step.} For each state $s$, update $C[s]$ to the mean of
  all training samples assigned to state $s$.
\end{enumerate}

Training data consists of i.i.d.\ samples from $\mathcal{N}(0, 1/128)$,
matching the post-FWHT distribution. This eliminates the need for model-specific
calibration data — the codebook depends only on the known post-rotation
distribution.

\textbf{Initialization.} Codebooks are initialized with analytically-computed
Lloyd-Max levels for the Gaussian distribution, arranged in monotonically
increasing order within each coset. This \emph{coset initialization} preserves
the set-partitioning structure that provides TCQ's coding gain.

\subsection{MSE--Quality Divergence}
\label{sec:mse_divergence}

A central finding of our work is that \textbf{minimizing quantization MSE does
not minimize downstream model quality loss}. Table~\ref{tab:mse_ppl} shows the
relationship between GLA iteration count, MSE reduction, and model perplexity.

\begin{table}[t]
\centering
\caption{GLA iteration sweep for 3-bit TCQ codebooks. MSE improves
monotonically; perplexity does not. Results on Qwen3.5-27B, wikitext-2,
64 chunks at 2K context.}
\label{tab:mse_ppl}
\begin{tabular}{@{}lcccc@{}}
\toprule
GLA Iters & MSE Reduction & PPL (2K) & $\Delta$ PPL \\
\midrule
0 (coset init) & 0.1\% & 5.919 & +1.64\% \\
3 & 24.3\% & 5.845 & +0.37\% \\
5 & 33.8\% & 5.858 & +0.58\% \\
10 & 40.4\% & 5.939 & +1.97\% \\
20 & 46.0\% & 5.971 & +2.53\% \\
30 & 48.2\% & 5.873 & +0.85\% \\
50 & 52.8\% & 5.831 & +0.13\% \\
100 & 54.1\% & 5.889 & +1.12\% \\
\bottomrule
\end{tabular}
\end{table}

Several patterns emerge:

\textbf{Early stopping is optimal.} Three GLA iterations suffice to capture
the trellis coding gain without degrading model quality. This result is robust
across seeds and datasets: we verified with 4 independent random seeds and 3
dataset splits (test, validation, train) with 64 chunks each, confirming that
3-iteration codebooks consistently match the best known codebook
($\Delta < 0.005$ PPL).

\textbf{Over-training degrades quality.} Beyond 10 iterations, perplexity
degrades by 0.06--0.12 despite continued MSE improvement. This degradation
persists across all evaluation datasets, ruling out measurement noise.

\textbf{Aggregate error metrics do not predict the degradation.} We tested 13
error statistics including covariance structure, kurtosis, tail percentiles,
signal-error correlation, and simulated attention divergence
(Table~\ref{tab:metrics_monotonic}). \emph{Every metric improves monotonically}
with GLA iterations, yet perplexity exhibits non-monotonic behavior. The
mechanism by which over-trained codebooks degrade model quality remains invisible
to all single-layer aggregate analysis we attempted.

\begin{table}[t]
\centering
\caption{All tested aggregate error metrics improve monotonically with GLA
  iterations, yet none predict perplexity. Pearson correlation with PPL shown.}
\label{tab:metrics_monotonic}
\begin{tabular}{@{}lcc@{}}
\toprule
Metric & Direction w/ training & Pearson w/ PPL \\
\midrule
Element-wise MSE & Improves & +0.10 \\
Off-diagonal $\Sigma_e$ ratio & Improves & +0.07 \\
Error kurtosis & Improves & +0.05 \\
Tail percentiles (p99.9) & Improves & +0.03 \\
$\mathrm{tr}(\Sigma_q \cdot \Sigma_e)$ & Improves & +0.05 \\
Signal-error correlation & Improves & +0.05 \\
Simulated attention KL & Improves & +0.05 \\
Temperature perturbation $\varepsilon$ & Improves & +0.05 \\
\bottomrule
\end{tabular}
\end{table}

\textbf{Connection to the quantization literature.} This finding has theoretical
grounding in three independent results: (1) Kasner \& Marcellin's Universal TCQ
demonstrated that untrained codebooks with uniform thresholds achieve
``performance comparable with fully optimized ECTCQ'' at most
rates~\cite{kasner1999utcq} — the trellis structure provides coding gain, not
codeword optimization; (2) Zamir \& Feder showed that lattice (structured)
quantizers produce white, signal-independent noise, while trained quantizers
produce signal-dependent noise that is catastrophically amplified by nonlinear
operations like softmax~\cite{zamir1996lattice}; (3) Ordentlich \& Polyanskiy
proved that for matrix multiplication (which attention is), MSE-optimal
quantizers are 3.4$\times$ worse than task-optimal ones at realistic
dimensions~\cite{ordentlich2024matmul}.

We hypothesize that over-training progressively destroys the coset regularity
that provides TCQ's coding gain. The coset structure — monotonically interleaved
reconstruction levels across trellis states — acts as an implicit regularizer
that maintains error whiteness. GLA iterations, optimizing only for MSE, have no
incentive to preserve this structure. We observe that monotonic coset fraction
degrades from 64/64 (coset init) to 48/64 after aggressive training.

\subsection{Context-Dependent Codebook Quality}
\label{sec:context_codebook}

Our systematic evaluation across context lengths from 2K to 128K tokens reveals
that \textbf{no single codebook is optimal at all context lengths}. Rankings
among codebooks --- including between trained and analytical (compiled-in)
codebooks, and between different training methods --- shift unpredictably with
context length. Codebook selection affects KLD by 2--7\%, a meaningful but
secondary effect compared to the 20--44\% gain from TCQ itself.

We propose a two-regime model that explains \emph{why} rankings are
context-dependent, grounded in the interaction between quantization error
statistics and the attention mechanism:

\subsubsection{Short Context ($< 8$K): Variance Dominates}

At short context lengths, softmax attention is concentrated on a small number of
tokens (low-entropy attention distribution). Individual quantization error
spikes can corrupt the attention pattern for an entire query. In this regime,
\emph{error variance} — the regularity and predictability of quantization
errors — dominates model quality.

Structured codebooks (fewer GLA iterations, preserved coset regularity) produce
more predictable errors. By the Zamir-Feder lattice noise
theorem~\cite{zamir1996lattice}, regular quantizer structures produce
approximately white, signal-independent noise. This noise is benign for softmax
attention because it does not systematically shift attention weights.

\subsubsection{Long Context ($> 8$K): Bias Dominates via CLT Averaging}

At long context lengths, softmax attention distributes weight across many
tokens. The output at each position is a weighted average of value vectors:

\begin{equation}
\mathbf{o} = \sum_{i=1}^{n} \alpha_i \mathbf{v}_i
\end{equation}

where $\alpha_i$ are the softmax attention weights. The quantization error in
the output is correspondingly a weighted average of per-token errors:

\begin{equation}
\mathbf{e}_o = \sum_{i=1}^{n} \alpha_i \mathbf{e}_i
\end{equation}

By the central limit theorem, as the effective number of attended tokens
$n_{\text{eff}} = 1 / \sum_i \alpha_i^2$ grows, the \emph{variance} of
$\mathbf{e}_o$ is suppressed by a factor of $\sim 1/n_{\text{eff}}$, while
the \emph{bias} (systematic error from suboptimal codebook placement) is not:

\begin{equation}
\Var[\mathbf{e}_o] \approx \frac{\sigma_e^2}{n_{\text{eff}}}, \qquad
\E[\mathbf{e}_o] = \sum_i \alpha_i \E[\mathbf{e}_i] \neq 0
\end{equation}

At long context, CLT averaging suppresses the variance component of output
error, changing which codebook properties matter. The relative importance of
MSE (favoring trained codebooks) versus error regularity (favoring analytical
codebooks) shifts with context length, producing the observed ranking
instability. At 3-bit, the analytical codebook's regularity advantage is
sufficient to win at moderate context lengths (8K--32K) despite higher MSE,
while trained codebooks may recover at very long contexts where even the
attenuated bias term matters.

\subsubsection{Supporting Theory: Finite Blocklength Dispersion}

The crossover is further explained by the finite-blocklength dispersion
framework of Kostina \& Verd\'{u}~\cite{kostina2012finite}. The gap between a
quantizer's distortion and the rate-distortion bound $R(D)$ decomposes as:

\begin{equation}
D - D^*(R) \approx \underbrace{D_{\text{mismatch}}}_{\text{constant}} +
\underbrace{\sqrt{\frac{V(R)}{n}} \cdot Q^{-1}(\epsilon)}_{\text{shrinks with } n}
\end{equation}

A well-trained codebook has small $D_{\text{mismatch}}$ and its gap is
dispersion-dominated (shrinks with block length). A mismatched codebook has
large, persistent $D_{\text{mismatch}}$. This framework predicts that codebook
differences become more consequential at long contexts, where the dispersion
term is small and codebook mismatch dominates --- consistent with our
observation that codebook rankings change with context length.

\subsubsection{The 2-Bit Regime: Overload Distortion}

At 2 bits per value, the codebook has only 4 reconstruction levels per coset.
In this regime, \emph{overload distortion} (clipping of tail samples) dominates
granular distortion. The UTCQ universality result — that untrained codebooks
suffice — breaks down because there are too few levels for the coset structure
alone to adequately cover the Gaussian tails.

This explains why codebook training provides substantially larger improvements
at 2-bit (up to 16\% KLD reduction) than at 3-bit, where the analytical
codebook is competitive at short contexts. It also explains the compiled-in 2-bit codebook's anomalous behavior
at 64K tokens: the specific GLA local minimum produces error tails (higher
kurtosis) that are catastrophically amplified by diluted softmax attention at
long context, while a different local minimum with similar MSE but lower
kurtosis handles long context gracefully.

\subsection{Query Anisotropy and Theoretical Error Ceiling}
\label{sec:anisotropy}

We measure the covariance structure of query vectors in the rotated
(post-FWHT) domain and find extreme anisotropy:

\begin{proposition}[Q Anisotropy]
For Qwen3.5-27B, the rotated query covariance $\Sigma_q$ has effective rank
$\approx 5$ out of 128 dimensions, with eigenvalue ratio
$\lambda_{\max}/\lambda_{\min} \approx 937$. The dominant principal component
captures $\sim 96\%$ of query variance.
\end{proposition}

This has a striking implication for quantization: the attention dot product
$\mathbf{q}^T \mathbf{k}$ is sensitive only to key errors along $\Sigma_q$'s
principal directions. The expected attention error is:

\begin{equation}
\E\left[\Var_j(\mathbf{q}^T \mathbf{e}_j)\right] =
\mathrm{tr}(\Sigma_q \cdot \Sigma_e)
\end{equation}

where $\Sigma_e$ is the quantization error covariance. If errors could be
concentrated in the 123 directions orthogonal to $\Sigma_q$'s principal
subspace, attention error would decrease by $\sim 25\times$ at the same MSE.

However, we find that TCQ quantization errors are \emph{isotropic} in
$\Sigma_q$'s eigenvector basis: the ratio of error variance along the top-5
query eigenvectors to average error variance is $1.00 \pm 0.01$ for all tested
codebooks. The FWHT rotation that enables analytical codebook design also
destroys any structure that could be exploited for query-aware quantization.

\textbf{Per-layer variation.} The anisotropy varies dramatically across layers:
layer 19 has effective rank 1.9 (nearly one-dimensional queries), while layer 20
has effective rank 114 (nearly isotropic). This suggests that per-layer
codebook optimization — weighting errors by layer-specific $\Sigma_q$ — could
yield significant quality improvements, but is architecturally
complex.

This finding connects to the Ordentlich-Polyanskiy result~\cite{ordentlich2024matmul}: the
optimal quantizer for the matrix multiplication $Q^T K$ has a fundamentally
different structure than the MSE-optimal quantizer, with density proportional to
$(f(x) \cdot \E[Y^2 | X = x])^{1/3}$ rather than $f(x)^{1/3}$. The 25$\times$
theoretical ceiling represents the gap between our current isotropic TCQ and a
hypothetical query-weighted quantizer.


% ============================================================
\section{Experiments}
\label{sec:results}

\subsection{Setup}
\label{sec:exp_setup}

\textbf{Model.} We evaluate on Qwen3.5-27B~(Q6\_K quantized weights), a
27-billion parameter dense transformer with 40 layers, 24 query heads, 4
key-value heads, and head dimension $d = 256$. The large head dimension is
favorable for rotation-based quantization: the post-FWHT distribution converges
more tightly to Gaussian at $d = 256$ than at $d = 128$, and each 256-element
vector provides a longer trellis path for TCQ to exploit.

\textbf{Hardware.} Primary measurements use an NVIDIA RTX 3090 (24\,GB VRAM,
Ampere SM86) for context lengths 2K--32K. Deep-context measurements (48K--128K)
use an NVIDIA A100 (80\,GB VRAM, Ampere SM80), as the KV cache exceeds 24\,GB
at these lengths. All measurements are single-GPU with full model offload.
Cross-GPU results are reported separately, as we observe that quantization
rankings are not stable across GPU architectures.

\textbf{Dataset.} WikiText-2 test set~(raw tokens, 1.29\,M bytes). We use
maximum chunks at each context length (8 chunks at $\leq$16K, 4 chunks at
24K--32K, 2 chunks at 48K, 1 chunk at 64K--128K) to minimize variance from
chunk sampling.

\textbf{Metrics.} Our primary metric is KL divergence between the output
logits of the quantized KV cache and an FP16 (lossless) KV cache baseline,
measured using \texttt{llama-perplexity --kl-divergence}. KLD directly measures
how much the quantization changes the model's output distribution, independent
of the downstream task. We also report perplexity (PPL) on WikiText-2 as a
secondary metric, noting that PPL and KLD can diverge --- a quantized model can
achieve lower PPL than FP16 if the quantization acts as a regularizer
(Section~\ref{sec:mse_divergence}).

\textbf{Base logits.} For each context length, we generate reference logits
using the same model with FP16 KV cache. Base logits are generated and measured
on the same GPU to avoid cross-platform artifacts.

\textbf{Baselines.} We compare six KV cache configurations:
\begin{itemize}\setlength{\itemsep}{0pt}\setlength{\parskip}{0pt}
  \item \textbf{f16}: Unquantized FP16 KV cache (reference).
  \item \textbf{q8\_0}: 8-bit scalar quantization (llama.cpp built-in).
  \item \textbf{turbo3}: 3-bit TurboQuant vector quantization without TCQ.
  \item \textbf{turbo3\_tcq}: 3-bit TurboQuant with trellis-coded quantization.
  \item \textbf{turbo2}: 2-bit TurboQuant without TCQ.
  \item \textbf{turbo2\_tcq}: 2-bit TurboQuant with TCQ.
\end{itemize}
This isolates the contribution of TCQ: turbo3 vs turbo3\_tcq and turbo2 vs
turbo2\_tcq differ \emph{only} in whether the trellis code is applied.

\textbf{Codebook configurations.} TCQ uses a 512-state bitshift trellis with
codebooks trained via the Generalized Lloyd Algorithm
(Section~\ref{sec:gla}). We evaluate both the analytically-derived
compiled-in codebook and product-aware trained codebooks, reporting whichever
achieves lower KLD at each context length. Non-TCQ types use the same FWHT
rotation and norm extraction pipeline but with scalar Lloyd-Max quantization.

\textbf{Alpha scaling.} The norm reconstruction scale $\alpha$ is optimized
per context length via grid search over $[0.96, 1.12]$, separately for keys
and values (Section~\ref{sec:norm}).

% ------------------------------------------------------------
\subsection{KLD Results}
\label{sec:kld_results}

Table~\ref{tab:kld_tcq} shows our primary result: TCQ consistently reduces KL
divergence by 10--44\% relative to scalar quantization at the same bit rate.

\begin{table}[t]
\centering
\caption{KL divergence (lower = better) for TCQ vs non-TCQ at all context
lengths. Compiled-in codebook, $\alpha_V = 1.04$.
RTX~3090 (2K--32K), A100 (48K--128K).}
\label{tab:kld_tcq}
\begin{tabular}{@{}rcccccc@{}}
\toprule
Context & turbo3 & turbo3\_tcq & $\Delta$ & turbo2 & turbo2\_tcq & $\Delta$ \\
\midrule
2K   & 0.0748 & \textbf{0.0552} & $-$26.1\% & 0.1363 & \textbf{0.1119} & $-$17.9\% \\
4K   & 0.0635 & \textbf{0.0571} & $-$10.2\% & 0.1401 & \textbf{0.1063} & $-$24.1\% \\
8K   & 0.0904 & \textbf{0.0746} & $-$17.4\% & 0.1797 & \textbf{0.1362} & $-$24.2\% \\
16K  & 0.0920 & \textbf{0.0702} & $-$23.7\% & 0.1862 & \textbf{0.1300} & $-$30.2\% \\
24K  & 0.0819 & \textbf{0.0686} & $-$16.3\% & 0.1672 & \textbf{0.1217} & $-$27.2\% \\
32K  & 0.0576 & \textbf{0.0447} & $-$22.3\% & 0.1383 & \textbf{0.0914} & $-$33.9\% \\
\midrule
48K  & 0.0669 & \textbf{0.0483} & $-$27.8\% & 0.1422 & \textbf{0.0953} & $-$33.0\% \\
64K  & 0.0699 & \textbf{0.0553} & $-$20.9\% & 0.1532 & \textbf{0.0999} & $-$34.8\% \\
128K & 0.0408 & \textbf{0.0306} & $-$25.0\% & 0.1091 & \textbf{0.0607} & $-$44.4\% \\
\bottomrule
\end{tabular}
\end{table}

Several patterns emerge from Table~\ref{tab:kld_tcq}:

\textbf{TCQ advantage is large and consistent.} At 3-bit, TCQ reduces KLD by
10--28\% across all context lengths. At 2-bit, the improvement is even larger:
18--44\%, reflecting the greater coding gain available at lower rates where
scalar quantization has fewer levels.

\textbf{2-bit improvement grows with context.} The 2-bit TCQ advantage
generally increases from 17.9\% at 2K to 44.4\% at 128K tokens (with
minor dips at 24K and 48K). This
context-scaling effect is explained by CLT averaging
(Section~\ref{sec:context_codebook}): at long contexts, the trellis code's
structured error pattern averages more favorably than scalar quantization's
independent errors.

\textbf{Non-monotonic context scaling.} KLD does not increase monotonically
with context length. Instead, there is a spike at 8K followed by a decline
to the lowest values at 32K and 128K. The 32K turbo3\_tcq KLD (0.0447) is
\emph{lower} than the 2K value (0.0552), confirming that CLT averaging in the
attention mechanism can more than compensate for the increased sequence length.
This pattern is consistent across all quantization types including q8\_0,
suggesting it is a property of the evaluation data rather than the quantization.

% ------------------------------------------------------------
\subsection{Perplexity Results}
\label{sec:ppl_results}

Table~\ref{tab:ppl} reports perplexity on WikiText-2 at short and deep
context lengths.

\begin{table}[t]
\centering
\caption{Perplexity (lower = better) on WikiText-2. RTX~3090 (2K, 32K),
A100 (48K--128K). turbo3\_tcq achieves \emph{lower} PPL than FP16 at every
context length.}
\label{tab:ppl}
\begin{tabular}{@{}lcccc@{}}
\toprule
Type & 2K & 32K & 48K & 128K \\
\midrule
f16         & 5.805 & 6.950 & 6.286 & 5.540 \\
q8\_0       & 5.839 & 6.951 & 6.271 & 5.526 \\
turbo3      & 5.850 & 7.088 & 6.346 & 5.573 \\
\textbf{turbo3\_tcq} & \textbf{5.777} & \textbf{6.862} & \textbf{6.178} & \textbf{5.461} \\
turbo2      & 6.079 & 7.492 & 6.835 & 6.042 \\
\textbf{turbo2\_tcq} & \textbf{6.005} & \textbf{7.099} & \textbf{6.470} & \textbf{5.609} \\
\bottomrule
\end{tabular}
\end{table}

\textbf{turbo3\_tcq beats lossless FP16.} At every context length measured,
3-bit TCQ achieves lower perplexity than uncompressed FP16 KV cache. This
counterintuitive result arises from the norm scaling factor $\alpha_V = 1.04$:
the slight upward scaling of value norms acts as an implicit regularizer,
sharpening the attention distribution and suppressing the contribution of
noisy low-weight tokens. The improvement grows with context: $-$0.47\% at 2K,
$-$1.26\% at 32K, $-$1.71\% at 48K, $-$1.42\% at 128K.

\textbf{2-bit TCQ closes the gap.} Without TCQ, 2-bit quantization incurs
a 4.7--9.1\% PPL penalty relative to FP16. With TCQ, this shrinks to
1.2--3.5\%. At 128K, turbo2\_tcq (5.609) is within 1.24\% of FP16 (5.540) ---
a practical operating point for memory-constrained deployments.

% ------------------------------------------------------------
\subsection{Codebook Selection Across Contexts}
\label{sec:codebook_context_results}

We trained codebooks using four variants of the Generalized Lloyd Algorithm:
\emph{vanilla} (standard GLA), \emph{mono} (monotonicity-constrained),
\emph{product} (product-aware distortion weighting), and \emph{product\_mono}
(both constraints). Each variant was trained for 100 iterations, producing 400
codebooks per bit width. Table~\ref{tab:codebook_3bit} shows the 3-bit context
scaling results for the top codebooks.

\begin{table}[t]
\centering
\caption{3-bit KLD across context lengths for top codebooks ($\alpha_V = 1.04$).
Bold = best at each context. The optimal codebook shifts with context length.}
\label{tab:codebook_3bit}
\begin{tabular}{@{}lcccccc@{}}
\toprule
Codebook & 2K & 4K & 8K & 16K & 24K & 32K \\
\midrule
compiled-in         & .0552 & .0571 & \textbf{.0746} & .0702 & .0686 & .0447 \\
pm/iter080          & \textbf{.0513} & .0567 & .0779 & .0697 & \textbf{.0630} & .0442 \\
pm/iter090          & .0595 & .0519 & .0788 & .0710 & .0638 & \textbf{.0419} \\
pm/iter060          & .0593 & \textbf{.0517} & .0793 & .0694 & .0663 & .0461 \\
\bottomrule
\end{tabular}
\end{table}

\textbf{No single codebook wins everywhere.} The compiled-in (analytical)
codebook wins at 8K; product\_mono/iter080 wins at 2K and 24K;
product\_mono/iter090 wins at 32K; product\_mono/iter060 wins at 4K. However,
the differences between codebooks are small (2--7\%) relative to the TCQ
vs non-TCQ gap (17--27\%). For deployment, product\_mono/iter080 achieves the
best average KLD (0.0605 vs 0.0617 for compiled-in, $-$2.0\%).

At 2-bit (Table~\ref{tab:codebook_2bit}), codebook training has a larger
effect, with the best trained codebook improving average KLD by 3.5\%.

\begin{table}[t]
\centering
\caption{2-bit KLD across context lengths for top codebooks ($\alpha_V = 1.04$).
Trained codebooks provide larger improvements than at 3-bit.}
\label{tab:codebook_2bit}
\begin{tabular}{@{}lcccccc@{}}
\toprule
Codebook & 2K & 4K & 8K & 16K & 24K & 32K \\
\midrule
compiled-in          & .1119 & .1063 & .1362 & \textbf{.1300} & .1217 & .0914 \\
pm/iter090           & \textbf{.0941} & .1016 & .1320 & .1352 & \textbf{.1200} & .0901 \\
vanilla/iter060      & .0951 & \textbf{.0957} & \textbf{.1283} & .1378 & .1279 & .0932 \\
pm/iter080           & .0964 & .1130 & .1307 & .1353 & .1278 & \textbf{.0894} \\
\bottomrule
\end{tabular}
\end{table}

\textbf{Product-aware training dominates.} The product\_mono training method
--- which jointly optimizes for the product distortion metric used in attention
while enforcing monotonic coset structure --- produces the best codebooks at
both bit widths. This is consistent with the Ordentlich-Polyanskiy finding that
MSE is suboptimal for matrix multiplication~\cite{ordentlich2024matmul}:
product-aware training steers codebooks toward attention-relevant distortion
rather than pure MSE.

\textbf{Convergence iteration differs by rate.} The optimal GLA iteration
is 80 for 3-bit and 90 for 2-bit. At 2-bit, more training iterations are
needed because the smaller codebook (4 levels per coset vs 8) has less
redundancy from the coset structure alone. This is consistent with the
breakdown of UTCQ universality at low rates
(Section~\ref{sec:context_codebook}).

% ------------------------------------------------------------
\subsection{Alpha Optimization}
\label{sec:alpha_results}

The norm scaling factor $\alpha_V$ (Section~\ref{sec:norm}) controls the
trade-off between quantization bias and variance. We sweep $\alpha_V \in
\{0.98, 1.00, 1.02, 1.04, 1.06, 1.08, 1.10, 1.12\}$ at each context length
for the selected codebooks. Table~\ref{tab:alpha} shows the optimal alpha and
the KLD improvement over the default $\alpha = 1.04$.

\begin{table}[t]
\centering
\caption{Optimal norm scaling $\alpha_V^*$ by context length.
3-bit alpha \emph{decreases} with context; 2-bit alpha \emph{increases} ---
opposite trends driven by different error regimes. $\Delta$ = improvement over
default $\alpha = 1.04$.}
\label{tab:alpha}
\begin{tabular}{@{}rcccccc@{}}
\toprule
& \multicolumn{3}{c}{3-bit (pm/iter080)} & \multicolumn{3}{c}{2-bit (pm/iter090)} \\
\cmidrule(lr){2-4}\cmidrule(lr){5-7}
Context & $\alpha^*$ & KLD & $\Delta$ & $\alpha^*$ & KLD & $\Delta$ \\
\midrule
2K  & 1.04 & 0.0513 & ---     & 1.04 & 0.0941 & --- \\
8K  & 1.02 & 0.0743 & $-$4.6\% & 1.08 & 0.1275 & $-$3.4\% \\
16K & 1.02 & 0.0651 & $-$6.6\% & 1.10 & 0.1255 & $-$7.1\% \\
32K & 1.00 & 0.0402 & $-$9.1\% & 1.10 & 0.0840 & $-$6.7\% \\
\midrule
48K$^\dagger$ & 1.02 & 0.0471 & $-$3.3\% & 1.08$^\dagger$ & 0.0877 & $-$3.5\% \\
64K$^\dagger$ & 1.00 & 0.0496 & $-$6.9\% & 1.08$^\dagger$ & 0.0926 & $-$7.3\% \\
\bottomrule
\multicolumn{7}{l}{\footnotesize $^\dagger$A100 (SM80); all others RTX~3090 (SM86). 2-bit 48K--64K uses vanilla/iter060.}
\end{tabular}
\end{table}

\textbf{Opposite alpha trends at 3-bit and 2-bit.} At 3-bit, the optimal alpha
decreases monotonically from 1.04 at 2K to 1.00 at 32K--64K. At 2-bit, the
opposite occurs: optimal alpha \emph{increases} from 1.04 at 2K to 1.08--1.10
at longer contexts. This asymmetry has a natural interpretation. At 3-bit (8
centroids per coset), the trellis code already provides good coverage of the
signal space, and the dominant error is bias from norm over-correction. Reducing
$\alpha$ toward 1.0 removes this bias. At 2-bit (4 centroids per coset), the
dominant error is variance from coarse quantization. Increasing $\alpha$ inflates
the reconstructed norms, which amplifies the signal relative to the quantization
noise --- effectively performing a form of noise-aware Wiener filtering.

\textbf{Alpha tuning is worth 3--9\% KLD.} The improvement from tuning alpha
grows with context length, from negligible at 2K to 6--9\% at 32K.

\textbf{Decode-time alpha is preferred for deployment.}
Alpha can be applied at encode time (baked into the fp16 norm before storage)
or at decode time (multiplied during dequantization). While both paths share
the same optimal alpha at most context lengths, decode-time application
produces 4\% lower KLD at 8K context for 3-bit and 3.5\% lower for 2-bit.
This advantage likely arises because decode-time scaling operates in float32
after the fp16 round-trip, avoiding interactions between the scaling factor
and fp16 rounding. More importantly, decode-time alpha enables
\emph{context-adaptive deployment}: the optimal $\alpha_V$ can be changed at
inference time based on the current context depth, without re-encoding any
cached keys or values. In our implementation, this is a single float
parameter passed to the dequantization kernel, controllable via environment
variable.

% ------------------------------------------------------------
\subsection{Deep Context (48K--128K)}
\label{sec:deep_context}

Context lengths beyond 32K exceed the KV cache capacity of FP16 on a 24~GB GPU.
Tables~\ref{tab:kld_tcq} and~\ref{tab:ppl} include measurements at 48K, 64K,
and 128K on an A100~80~GB GPU with separate same-GPU base logits (these numbers
are not directly comparable to RTX~3090 measurements due to different float
rounding). Deep context measurements use fewer chunks (2 at 48K, 1 at
64K--128K vs 8 at $\leq$16K) and consequently have higher variance.

Two findings extend from the shorter-context results:

\textbf{TCQ advantage grows at deep context.} The 2-bit TCQ improvement
increases from 17.9\% at 2K to 44.4\% at 128K (Table~\ref{tab:kld_tcq}). At
128K, turbo2\_tcq achieves KLD 0.0607 --- within 2$\times$ of turbo3 without
TCQ (0.0408). The trellis code's structured errors average more favorably under
the CLT effect in attention (Section~\ref{sec:context_codebook}), and this
benefit compounds at longer sequences.

\textbf{128K context on 24~GB.} With turbo3 at 3.5~bpv, the KV cache for 128K
tokens of Qwen3.5-27B occupies $\sim$4.7~GB, leaving ample headroom on a 24~GB
GPU alongside the 20~GB model. FP16 KV cache would require $\sim$21~GB at 128K,
exceeding the GPU budget. turbo2\_tcq at 2.25~bpv reduces this further to
$\sim$3.0~GB. In our competitive evaluation
(Section~\ref{sec:discussion}), all four third-party TurboQuant forks
failed at 32K on 24~GB --- only our implementation succeeded.

% ------------------------------------------------------------
\subsection{Decode Speed}
\label{sec:speed}

TCQ is not free. The sliding-window trellis decode
(Section~\ref{sec:decode}) replaces a direct 3-bit table lookup with a
9-bit state-indexed extraction per group of three values. Table~\ref{tab:speed}
reports decode throughput at four context depths.

\begin{table}[t]
\centering
\caption{Decode throughput (tok/s, higher = better) on RTX~3090 at
increasing KV cache depth. \texttt{tg128}, 3 repetitions.}
\label{tab:speed}
\begin{tabular}{@{}lcccc@{}}
\toprule
Type & 2K & 8K & 16K & 32K \\
\midrule
turbo3       & 28.24 & 26.14 & 24.32 & 21.72 \\
turbo3\_tcq  & 26.99 & 24.01 & 21.43 & 17.69 \\
\quad \emph{overhead} & \emph{$-$4.4\%} & \emph{$-$8.1\%} & \emph{$-$11.9\%} & \emph{$-$18.6\%} \\
\midrule
turbo2       & 28.30 & 26.49 & 24.88 & 22.31 \\
turbo2\_tcq  & 27.35 & 24.57 & 21.85 & 17.97 \\
\quad \emph{overhead} & \emph{$-$3.4\%} & \emph{$-$7.2\%} & \emph{$-$12.2\%} & \emph{$-$19.4\%} \\
\bottomrule
\end{tabular}
\end{table}

\textbf{TCQ overhead grows with context depth.} At 2K, the sliding-window decode
adds only 3--4\% overhead. At 32K, this grows to 18--19\%. The reason is
architectural: as the KV cache grows, the attention inner loop --- where the
trellis decode operates --- occupies a larger fraction of total compute. At 2K,
attention is a minor cost and the FFN layers dominate; at 32K, the KV cache
read becomes the bottleneck and the trellis decode overhead is fully exposed.

\textbf{Memory efficiency.} turbo3\_tcq at 3.25~bpv uses 7\% less memory than
turbo3 at 3.5~bpv, and 59\% less than q8\_0 at 8~bpv. turbo2\_tcq at 2.25~bpv
provides 72\% memory savings over q8\_0. For a 27B model at 32K context, this is
the difference between fitting on a 24~GB consumer GPU and requiring 48~GB+
server hardware.

% ============================================================
\section{Discussion and Conclusion}
\label{sec:discussion}

We have presented the first application of trellis-coded quantization to the KV
cache of large language models. Our approach augments the existing TurboQuant
vector quantization pipeline --- FWHT rotation, Lloyd-Max scalar quantization,
norm extraction --- with a 512-state bitshift trellis that encodes each quantized
value as a state transition rather than an independent index. The key results:

\begin{enumerate}\setlength{\itemsep}{2pt}
\item \textbf{17--44\% KLD reduction} over scalar quantization at the same bit
rate, across all context lengths from 2K to 128K (Table~\ref{tab:kld_tcq}).
\item \textbf{turbo3\_tcq beats FP16 in perplexity} at every context length
measured (Table~\ref{tab:ppl}), achieving 3.5~bpv with better quality than
lossless 16-bit storage.
\item \textbf{Norm alpha scaling} with opposite optimal trends at 3-bit
(decreasing with context) and 2-bit (increasing), yielding 3--9\% KLD
improvement (Table~\ref{tab:alpha}). Decode-time alpha application is
preferred: it matches or exceeds encode-time quality and enables
context-adaptive deployment without re-encoding the cache.
\item \textbf{O(1) decode} via sliding-window trellis extraction, with 4--19\%
throughput overhead depending on context depth (Table~\ref{tab:speed}).
\item \textbf{128K context on 24~GB} consumer GPUs --- a capability no 8-bit
scheme can match on this model.
\end{enumerate}

\textbf{Why does quantization beat lossless?} The turbo3\_tcq PPL result (5.777
vs FP16's 5.805 at 2K) is surprising but has a clean explanation. The norm
scaling factor $\alpha_V$ (1.04 at 2K context) slightly inflates value norms, which sharpens
the attention distribution by amplifying the contribution of high-weight tokens
relative to low-weight noise tokens. This is equivalent to a temperature
reduction in the softmax. The trellis code further improves this by making
quantization errors correlated across adjacent values, which creates a smoother
error surface that the attention softmax can more easily absorb. The effect is
a form of implicit regularization: the quantized model attends more precisely
than the unquantized model, at the cost of a small distributional shift (KLD
0.055 at 2K).

\textbf{Codebook universality.} Our experiments confirm the theoretical
prediction of Section~\ref{sec:context_codebook}: after FWHT rotation, the
signal is sufficiently Gaussian that trained codebooks provide only 2--3.5\%
improvement over analytical codebooks at 3-bit. At 2-bit, where the Gaussianity
assumption breaks down (fewer centroids = more sensitivity to tail behavior),
training provides more benefit, and the product-aware distortion metric outperforms
pure MSE --- consistent with the Ordentlich-Polyanskiy finding that MSE is
suboptimal for matrix multiplication~\cite{ordentlich2024matmul}.

\textbf{Competitive landscape.} We benchmarked four independent TurboQuant forks
(TheTom, Madreag, AmesianX, Duster) on the same GPU with the same base logits.
None implement TCQ. Our TCQ achieves 17--34\% lower KLD than the best competitor
at every context length. At 32K, all four competitors failed while our
implementation succeeded. The speed gap is modest: the fastest competitor
(Madreag) achieves $\sim$5\% higher throughput than our non-TCQ baseline, likely
due to a larger block size (128 vs 32 elements). Our TCQ overhead (8\% at 8K)
is well within this margin.

\textbf{Limitations.} The TCQ decode overhead grows to 19\% at 32K context,
which may matter for latency-sensitive applications at long context. Prefill is
also slower: the Viterbi encode path (512 states $\times$ 128 forward steps)
adds $\sim$21\% overhead to prefill throughput. Our current implementation
supports only GQA architectures with $\text{head\_dim} = 128$ or $256$; models
with sliding-window attention (e.g., Gemma-3) require modifications to the V
un-rotation path.

\textbf{Future work.} Three directions are immediately actionable. First,
\emph{asymmetric K/V quantization}: using 3-bit TCQ for keys (which the softmax
exponentially amplifies) and 2-bit TCQ for values (which enter linearly) would
achieve an effective 2.75~bpv. Preliminary PPL results are mixed --- the
optimal K/V rate allocation likely depends on context length and warrants
systematic investigation. Second,
\emph{per-layer alpha or rate allocation}: our alpha is uniform
across layers, but per-layer or per-head schedules may yield further gains;
similarly, early layers may tolerate coarser quantization than later layers,
enabling mixed 2-bit/3-bit schedules. Third, \emph{automatic alpha selection}:
our decode-time alpha implementation already enables context-adaptive scaling
without re-encoding the KV cache; a closed-form mapping from context depth to
optimal $\alpha_V$ would eliminate manual tuning entirely.

% ============================================================
\section*{Disclosure}

This paper was written collaboratively by a human researcher and Claude, an AI
assistant by Anthropic. Claude contributed to literature review, theoretical
framing, drafting, and code review throughout the project. The research was
vibed, but the benchmarks weren't --- all experimental results were produced by
real CUDA kernels running on real GPUs with no AI in the measurement loop.

% ============================================================
% References — will convert to .bib later
\begin{thebibliography}{99}

\bibitem{turboquant2026}
TurboQuant: Online Vector Quantization with Near-Optimal Distortion Rate.
Google Research \& NYU. ICLR 2026. arXiv:2504.19874.

\bibitem{marcellin1990tcq}
M.~W. Marcellin and T.~R. Fischer.
Trellis coded quantization of memoryless and Gauss-Markov sources.
\textit{IEEE Trans. Comm.}, 38(1):82--93, 1990.

\bibitem{fischer1992ectcq}
T.~R. Fischer and M.~Wang.
Entropy-constrained trellis coded quantization.
\textit{IEEE Trans. Inform. Theory}, 38(2):415--426, 1992.

\bibitem{kasner1999utcq}
J.~H. Kasner and M.~W. Marcellin.
Universal trellis coded quantization.
\textit{IEEE Trans. Image Processing}, 8(12):1677--1687, 1999.

\bibitem{zamir1996lattice}
R.~Zamir and M.~Feder.
On lattice quantization noise.
\textit{IEEE Trans. Inform. Theory}, 42(4):1152--1159, 1996.

\bibitem{ordentlich2024matmul}
O.~Ordentlich and Y.~Polyanskiy.
Optimal quantization for matrix multiplication.
arXiv:2410.13780, 2024.

\bibitem{asymkv2024}
AsymKV: Asymmetric quantization of key and value for KV cache compression.
arXiv:2410.13212, 2024.

\bibitem{sabin1986gla}
M.~J. Sabin and R.~M. Gray.
Global convergence and empirical consistency of the generalized Lloyd algorithm.
\textit{IEEE Trans. Inform. Theory}, 32(2):148--155, 1986.

\bibitem{kostina2012finite}
V.~Kostina and S.~Verd\'{u}.
Fixed-length lossy compression in the finite blocklength regime.
\textit{IEEE Trans. Inform. Theory}, 58(6):3309--3338, 2012.

\bibitem{qtip2024}
A.~Tseng et al.
QTIP: Quantization with Trellises and Incoherence Processing.
NeurIPS 2024. arXiv:2406.11235.

\bibitem{quipsharp2024}
J.~Chee et al.
QuIP\#: Even Better LLM Quantization with Hadamard Incoherence and Lattice
Codebooks. arXiv:2402.04396, 2024.

\bibitem{gersho1979}
A.~Gersho.
Asymptotically optimal block quantization.
\textit{IEEE Trans. Inform. Theory}, 25(4):373--380, 1979.

\bibitem{gray1998quantization}
R.~M. Gray and D.~L. Neuhoff.
Quantization.
\textit{IEEE Trans. Inform. Theory}, 44(6):2325--2383, 1998.

% === Section 2 references ===

\bibitem{llama_cpp}
G.~Gerganov et al.
llama.cpp: LLM inference in C/C++.
{\small\url{https://github.com/ggerganov/llama.cpp}}, 2023--2026.

\bibitem{kivi2024}
Z.~Liu et al.
KIVI: A tuning-free asymmetric 2bit quantization for KV cache.
arXiv:2402.02750, 2024.

\bibitem{kvquant2024}
C.~Hooper et al.
KVQuant: Towards 10 million context length LLM inference with KV cache
quantization. arXiv:2401.18079, 2024.

\bibitem{vecinfer2025}
VecInfer: Outlier-suppressed vector quantization for memory-efficient LLM
inference. arXiv:2510.06175, 2025.

\bibitem{nsnquant2025}
NSNQuant: Calibration-free VQ for KV cache via triple normalization.
arXiv:2505.18231, 2025.

\bibitem{kvtc2025}
KVTC: KV cache transform coding.
NVIDIA Research. ICLR 2026. arXiv:2511.01815.

\bibitem{morekeyslessvalues2025}
More keys, less values: Asymmetric KV cache bit allocation.
arXiv:2502.15075, 2025.

\bibitem{rotatekv2025}
RotateKV: 2-bit KV cache quantization via outlier-aware adaptive rotations.
arXiv:2501.16383, 2025.

\bibitem{mixquant2026}
MixQuant: Pushing the limits of block rotations in post-training quantization.
arXiv:2601.22347, 2026.

\bibitem{ailon2009fast}
N.~Ailon and B.~Chazelle.
The fast Johnson-Lindenstrauss transform and approximate nearest neighbors.
\textit{SIAM J. Comput.}, 39(1):302--322, 2009.

\bibitem{lookabaugh1989}
T.~Lookabaugh and R.~M. Gray.
High-resolution quantization theory and the vector quantizer advantage.
\textit{IEEE Trans. Inform. Theory}, 35(5):1020--1033, 1989.

\end{thebibliography}

\end{document}
